%!TEX TS-program = xelatex
\documentclass{ausgabeOTH}
\usepackage[utf8]{inputenc}
\usepackage[german]{babel}
\usepackage[doublespacing]{setspace}
\usepackage{tabularx}
\makeatother
%-----Allgemeine Variablen----%

%Masterarbeit oder Bachelorarbeit
\newcommand\MaBaArbeit{Bachelorarbeit}
%Name des Studenten Eintragen(Name, Vorname)
\newcommand\Student{Superfleissiger Student}
%Studiengang des Studenten eintragen
\newcommand\Studiengang{Industrie 4.0 - Informatik}
%Matrikelnummer
\newcommand\Matrikelnummer{123445678}
%Semester eintragen
\newcommand\Semester{8}
%Ausgabedatum der Arbeit
\newcommand\Ausgabe{01.01.2020}
%Festgelegetes Abgabedatum
\newcommand\Abgabe{01.01.2025}
%Titel der Arbeit eintragen
\newcommand\Titel{Megatoller Titel in deutsch}
\newcommand\TitelEnglisch{Super awsome title in english}
%Kurzbeschreibung
\newcommand\Beschreibung{Total tolle und aussagekräftige Beschreibung der megaspannenden Aufgabenstellung der Abschlussarbeit mit dem super genialen Titel}
%Name des Erstprüfers
\newcommand\Aufgabensteller{Prof. Dr. Schlaubi Schlumpf}
%Name des Zweitprüfers
\newcommand\Zweitpruefer{Prof. Dr. Gargamel}



%----externe Bachelorarbeit----%
%Name der Firma
\newcommand\Firma{FIRMA}
%Firmenanschrift
\newcommand\Anschrift{ANSCHRIFT}
%Name des Betreuers innerhalb der Firma
\newcommand\Firmenbetreuer{NAME}
%Abteilung der Firma
\newcommand\Abteilung{ABTEILUNG}
%Telefonnummer
\newcommand\Telefon{12345678}



%--Änderungen--
%Zeitraum um den die Abgabe verlängert werden soll in Zahlen
\newcommand\Zeitraum{Wochen}
%Neuer Abgabetermin
\newcommand\AbgabeNeu{01.01.3000}
%Gründe für die Verzögerung
\newcommand\Grund{akzeptable Gründe sind z.B.: unvorhergesehene Komplexität des Themas,schwierigkeiten bei der Recherche(keine Literatur vorhanden), lange nachgewiesene Krankheit, Pflege von Familienangehörigen(nicht Kinderbetreuung)}





\makeatletter


\begin{document}


\subsection*{Ausgabe des Themas der \MaBaArbeit}
\vspace*{-.5cm}
\noindent\rule[1ex]{\textwidth}{1pt} %für durchgezogene Linie
%\vspace*{-.5cm}
\small{(zu § 12 der Allgemeinen Prüfungsordnung der Ostbayerischen Technischen Hochschule Amberg-Weiden)}
\normalsize



Studiengang: \Studiengang \hfill Semester: \Semester\\
Kandidat/Kandidatin: der Bachelorarbeit \hfill \Student \\
Aufgabensteller/Kandidat/in der Bachelorarbeit:\hfill \Aufgabensteller\\
Zweitprüfer:\hfill \Zweitpruefer\\

Die Bachelorarbeit wird außerhalb der Hochschule durchgeführt bei

Firma/Einrichtung:\hfill \Firma\\
Betreuer in der Firma/Einrichtung: \hfill \Firmenbetreuer\\
Telefon:\hfill \Telefon\\

Titel der Bachelorarbeit\\
Deutscher Titel: \Titel

Englischer Titel: \TitelEnglisch

Beschreibung der Aufgabenstellung: \Beschreibung

\vfill


\textbf{Ausgabedatum: \Ausgabe}\hfill \textbf{Abgabetermin: \Abgabe}\\
\vspace*{1cm}

\rule{0.3\textwidth}{1pt}\hfill \rule{0.3\textwidth}{.5pt}\\
\tiny{Kandidat/in (Unterschrift)} \hfill \tiny{Betreuer/in (Unterschrift)}\\



\end{document}
